% Generated by tex/build_swebok_structure.py from tex/chapters/ch01/09_実務上の考慮.tex
\subsubsection{要求の優先順位付け}

\term{要求の優先順位付け}は、限られた費用と時間の中で、価値の高い要求から扱うために必要である。優先度は、価値だけでなく、未実現だった場合の不満、実装費用、保守費用、技術リスク、利用されないリスクなどを考慮して決める。

\MiniBox{Kanoモデルの考え方}{ある機能があると満足度が上がる場合と、ないと強い不満が出る場合は異なる。たとえばメールソフトでは、迷惑メールフィルタはあると嬉しい機能かもしれないが、添付ファイルを扱えないと利用者は強く不満を持つ。優先度は満足と不満の両方から考える。}

優先度の表し方には、「必須」「望ましい」「あればよい」のような段階、1から10の数値、優先順位リストなどがある。実務では、細かな数値差よりも、同程度の優先度を持つ要求群を見つけることが重要である。
